\section{The finite prime-target operator}
\label{sec:local-operator}

Throughout this section $q\geq3$ is odd, $(q,h)=1$, and
$G_q=(\Z/q\Z)^\times$.  We use normalized counting measure on $G_q$.
For $\beta:G_q\to\C$, define
\begin{align}
 (\mathsf P_{q,h}\beta)(a)
 &=\frac1{\varphi(q)}
   \sum_{\substack{b\in G_q\\(ab+h,q)=1}}\beta(b),
 \label{eq:P-definition}\\
 \kappa_q
 &=\prod_{p\mid q}\frac{p-2}{p-1},
 \qquad
 \mathsf T_{q,h}=\kappa_q^{-1}\mathsf P_{q,h}.
 \label{eq:T-definition}
\end{align}
For each $a$, exactly
\begin{equation}\label{eq:row-degree}
 R_q=\varphi(q)\kappa_q
 =\prod_{p^e\parallel q}p^{e-1}(p-2)
\end{equation}
values of $b$ are admissible.  Consequently $\mathsf T_{q,h}$ is a Markov
operator.  Its kernel is real and symmetric, because $ab+h=ba+h$, so it is
self-adjoint on $L^2(G_q)$.

\subsection{Inclusion-exclusion and the arithmetic involutions}

For every squarefree $d\mid\rad(q)$ let
\begin{equation}\label{eq:J-definition}
 J_{h,d}(a)=-h a^{-1}\pmod d,
 \qquad a\in G_d,
\end{equation}
and let $U_{h,d}f=f\circ J_{h,d}$.  Then $J_{h,d}^2$ is the identity, so
$U_{h,d}$ is a self-adjoint unitary involution.  Define the normalized fiber
average
\begin{equation}\label{eq:fiber-average}
 (E_{q,d}\beta)(u)
 =\frac{\varphi(d)}{\varphi(q)}
  \sum_{\substack{b\in G_q\\b\equiv u\pmod d}}\beta(b).
\end{equation}

\begin{proposition}[Exact involution bundle]
\label{prop:involution-bundle}
For every $\beta:G_q\to\C$ and $a\in G_q$,
\begin{equation}\label{eq:involution-bundle}
 (\mathsf P_{q,h}\beta)(a)
 =\sum_{d\mid\rad(q)}
   \frac{\mu(d)}{\varphi(d)}
   \bigl(U_{h,d}E_{q,d}\beta\bigr)(a\bmod d).
\end{equation}
All terms in this identity are finite and exact.
\end{proposition}

\begin{proof}
Inclusion-exclusion gives
\[
 \one_{(ab+h,q)=1}
 =\sum_{\substack{d\mid\rad(q)\\ab\equiv-h\pmod d}}\mu(d).
\]
For fixed $a$ and $d$, the congruence selects the fiber
$b\equiv J_{h,d}(a)\pmod d$.  By \eqref{eq:fiber-average}, its unnormalized
sum is
\[
 \frac{\varphi(q)}{\varphi(d)}
 (E_{q,d}\beta)(J_{h,d}(a)).
\]
Divide by $\varphi(q)$ and sum over $d$.
\end{proof}

This formula identifies the exact link with the target-divisor involutions
of \citet{WangTPCSix2026}.  It is not a mixing assertion: each $U_{h,d}$ is
an involution, while the damping comes from the normalized combination of
all divisor layers and from the fiber averages.

\subsection{Complete character-block decomposition}

Let $\wh G_q$ be the character group.  We use the Fourier convention
\begin{equation}\label{eq:Fourier-convention}
 B(\chi)=\frac1{\varphi(q)}\sum_{b\in G_q}\beta(b)\chi(b),
 \qquad
 \beta(b)=\sum_{\chi\in\wh G_q}B(\chi)\overline{\chi(b)}.
\end{equation}
Write $q=\prod_{p^e\parallel q}p^e$.  Let $f_p(\chi)$ be the conductor
exponent of the local component $\chi_p$ modulo $p^e$.

Call $\chi$ \emph{surviving} if $f_p(\chi)\leq1$ for every $p\mid q$.
For such a character put
\begin{equation}\label{eq:exact-squarefree-conductor}
 r_\chi=\prod_{\substack{p\mid q\\f_p(\chi)=1}}p,
 \qquad
 d(r_\chi)=\prod_{p\mid r_\chi}\frac1{p-2},
\end{equation}
and
\begin{equation}\label{eq:c-multiplier}
 c_h(\chi)
 =(-1)^{\omega(r_\chi)}d(r_\chi)
   \overline{\chi(-h)}.
\end{equation}
The principal character is surviving, with $r_\chi=1$ and $c_h(\chi)=1$.

\begin{theorem}[Exact local spectrum]
\label{thm:exact-local-spectrum}
For every $\beta:G_q\to\C$,
\begin{equation}\label{eq:P-Fourier-formula}
 (\mathsf P_{q,h}\beta)(a)
 =\kappa_q
  \sum_{\substack{\chi\in\wh G_q\\\chi\ \mathrm{surviving}}}
  B(\chi)c_h(\chi)\chi(a).
\end{equation}
Equivalently,
\begin{equation}\label{eq:T-on-character}
 \mathsf T_{q,h}\overline\chi
 =\begin{cases}
   c_h(\chi)\chi,&\chi\text{ surviving},\\
   0,&\chi\text{ not surviving}.
  \end{cases}
\end{equation}
In particular, on a surviving mode,
\begin{equation}\label{eq:T-square-character}
 \mathsf T_{q,h}^{2}\chi=d(r_\chi)^2\chi.
\end{equation}

The nonzero spectrum consists of $1$, the signed eigenvalues on real
characters, and the pairs $\pm d(r_\chi)$ on nonreal conjugate pairs.  Zero
has multiplicity
\begin{equation}\label{eq:zero-multiplicity}
 \varphi(q)-\varphi(\rad(q)).
\end{equation}
For every integer $k\geq1$,
\begin{equation}\label{eq:even-trace}
 \tr(\mathsf T_{q,h}^{2k})
 =\prod_{p\mid q}\left(1+(p-2)^{1-2k}\right).
\end{equation}
\end{theorem}

\begin{proof}
It is enough to compute locally.  Fix $p^e\parallel q$, write $a$ for a
unit modulo $p^e$, and consider the sum of $\overline{\chi_p(b)}$ over
unit $b$ with $ab+h\not\equiv0\pmod p$.

If $f_p=0$, the summand is $1$ and the sum is
$p^{e-1}(p-2)$.  If $f_p=1$, the complete unit sum is zero and the excluded
fiber has $p^{e-1}$ points, on all of which the character has the same
value.  The allowed sum is therefore
\[
 -p^{e-1}\overline{\chi_p(-h)}\chi_p(a).
\]
If $f_p\geq2$, both the complete unit sum and the sum over the excluded
fiber vanish: the restriction of $\chi_p$ to the kernel of reduction modulo
$p$ is nonprincipal.  Hence this local factor is zero.

Multiplying the local sums and dividing by $\varphi(q)$ proves
\eqref{eq:P-Fourier-formula} and \eqref{eq:T-on-character}.  Since
$c_h(\overline\chi)=\overline{c_h(\chi)}$ and
$|c_h(\chi)|=d(r_\chi)$, applying $\mathsf T_{q,h}$ twice proves
\eqref{eq:T-square-character}.  There are $p-1$ local characters with
conductor exponent at most one, so the nonzero subspace has dimension
$\prod_{p\mid q}(p-1)=\varphi(\rad(q))$.  Finally, group the surviving
characters by exact conductor.  At a prime $p$, there is one principal
choice and $p-2$ nonprincipal choices.  Thus
\[
 \sum_{\chi\ \mathrm{surviving}}d(r_\chi)^{2k}
 =\prod_{p\mid q}\left(1+(p-2)(p-2)^{-2k}\right),
\]
which is \eqref{eq:even-trace}.
\end{proof}

\begin{corollary}[Gap and the prime $3$ obstruction]
\label{cor:local-gap}
Suppose $q$ is squarefree.  If $3\mid q$, a nonprincipal conductor-$3$
mode has modulus-one eigenvalue, so there is no spectral gap on the
mean-zero subspace.  If $3\nmid q$ and $p_{\min}$ is the least prime
dividing $q$, then
\begin{equation}\label{eq:mean-zero-norm}
 \|\mathsf T_{q,h}\|_{\one^\perp\to\one^\perp}
 =\frac1{p_{\min}-2}<1.
\end{equation}
For nonsquarefree $q$, the same norm formula holds on $\one^\perp$;
equivalently, it holds on the mean-zero part of the surviving subspace.
The orthogonal complement of the surviving subspace is killed in one step.
\end{corollary}

The finite chain can therefore have a genuine local gap.  This does not say
that the primes, or the sequence $\Lambda(mn+h)$, form a mixing dynamical
system.  The distinction will be maintained throughout.
